% =========================================================================
% Capítulo 1: Cinemática
% =========================================================================
% Archivo: C1_00_introduccion.tex
% Descripción: Movimiento de puntos y sistemas de referencias
% =========================================================================

\chapter{Cinemática}

\section*{Introducción}

La cinemática es la rama de la mecánica que describe el movimiento sin considerar las causas que lo producen (fuerzas). En este capítulo estudiaremos cómo representar la posición, velocidad y aceleración de puntos materiales en diferentes sistemas de coordenadas, y cómo derivar vectores expresados en marcos de referencia que se mueven o rotan entre sí. La herramienta central será el Teorema del Transporte, que usaremos de manera sistemática en todos los problemas del libro.

\subsection*{Sistemas de Coordenadas}

El análisis cinemático depende críticamente de la elección del sistema de coordenadas. En este libro estudiaremos:

\begin{enumerate}
    \item \textbf{Coordenadas Cartesianas}: El sistema más simple e intuitivo.
    \item \textbf{Coordenadas Cilíndricas}: Naturales para problemas con simetría axial.
    \item \textbf{Coordenadas Esféricas}: Útiles para sistemas con simetría central.
    \item \textbf{Coordenadas Intrínsecas}: Basadas en la geometría de la trayectoria.
\end{enumerate}

La Figura \ref{fig:c1_intro_coordenadas_esfericas} muestra el sistema de coordenadas esféricas con sus vectores unitarios $\{\hat{e}_r, \hat{e}_\theta, \hat{e}_\phi\}$, que junto con las coordenadas cilíndricas será uno de los sistemas más usados en los problemas de este capítulo.

% =========================================================================
% Figura: Sistema de coordenadas esféricas
% =========================================================================
% Archivo: Coordenadas_esfericas.tikz.tex
% Usada en: C1_00_introduccion.tex
% Convención del libro: theta = ángulo azimutal (plano xy), phi = ángulo polar
% =========================================================================

\begin{figure}[!ht]
    \centering
    \tdplotsetmaincoords{70}{110}
    \begin{tikzpicture}[tdplot_main_coords, >=Latex,
                        ejes/.style={->, thick, gris1},
                        univec/.style={->, very thick, azul},
                        posvec/.style={->, very thick, myred}]

        % --- CONFIGURACIÓN DE PARÁMETROS ---
        \pgfmathsetmacro{\phivec}{48.17}   % Ángulo polar [deg]
        \pgfmathsetmacro{\thetavec}{63.5}  % Ángulo azimutal [deg]
        \def\rlen{6}                       % Largo del vector posición [u]

        % --- SISTEMAS DE COORDENADAS ---
        \draw[ejes] (0,0,0) -- (6.5,0,0) node[pos=1.1] {$x$};
        \draw[ejes] (0,0,0) -- (0,6,0) node[pos=1.05] {$y$};
        \draw[ejes] (0,0,0) -- (0,0,5.5) node[pos=1.05] {$z$};

        % --- PUNTOS Y COORDENADAS ---
        \tdplotsetcoord{P}{\rlen}{\phivec}{\thetavec}

        % --- VECTORES Y FUERZAS ---
        % Vector posición
        \draw[posvec] (0,0,0) -- (P) node[midway, above left] {$\vec{r}$};
        % Vectores unitarios en el punto P
        \tdplotsetcoord{Er}{7}{\phivec}{\thetavec}
        \draw[univec] (0,0,0) -- (Er) node[pos=1.05] {$\hat{e}_r$};
        \tdplotsetcoord{Et}{1}{90}{90+\thetavec}
        \draw[univec] (Pxy) -- ($(Et) + (Pxy)$) node[pos=1.3] {$\hat{e}_\theta$};
        \tdplotsetcoord{Ep}{1}{90+\phivec}{\thetavec}
        \draw[univec] (P) -- ($(Ep) + (P)$) node[pos=1.3] {$\hat{e}_\phi$};

        % --- ANOTACIONES Y ETIQUETAS ---
        % Líneas de ayuda (proyecciones)
        \draw[thick] (0,0,0) -- (Pxy);
        \draw[mydashed, gris1!60] (P) -- (Pxy);
        \draw[mydashed, gris1!60] (Px) -- (Pxy) node[pos=-0.15] {$x$};
        \draw[mydashed, gris1!60] (Py) -- (Pxy) node[pos=-0.3] {$y$};
        \draw[mydashed, gris1!60] (Pz) -- (P) node[pos=-0.15] {$z$};
        % Punto P
        \node[fill=gris1, circle, inner sep=0.9pt, label={above right:$P$}] at (P) {};
        % Ángulos
        \tdplotdrawarc{(0,0,0)}{0.7}{0}{\thetavec}{below}{$\theta$}
        \tdplotsetthetaplanecoords{\thetavec}
        \tdplotdrawarc[tdplot_rotated_coords]{(0,0,0)}{0.5}{0}{\phivec}{}{}
        \node at (0,0.25,0.67) {$\phi$};

    \end{tikzpicture}
    \caption{Coordenadas esféricas y sus vectores unitarios en el punto $P$}
    \label{fig:c1_intro_coordenadas_esfericas}
\end{figure}


\subsection*{Marcos de Referencia}

Un concepto fundamental es la distinción entre:
\begin{description}
    \item[Marcos Inerciales:] Sistemas de referencia en los que las leyes de Newton se cumplen sin términos de corrección.
    \item[Marcos No-Inerciales:] Sistemas acelerados o en rotación donde aparecen fuerzas ficticias.
\end{description}

\subsection*{Derivadas en Marcos Móviles}

El \textbf{Teorema del Transporte} permite relacionar las derivadas de un vector en dos marcos de referencia que rotan uno respecto al otro:

\[\transport{\vec{a}}{A}{B}\]

donde $\mathcal{A}$ es el marco fijo (inercial), $\mathcal{B}$ es el marco móvil y $\velang{A}{B}$ es la velocidad angular del marco $\mathcal{B}$ respecto al marco $\mathcal{A}$. Esta es la misma notación que usaremos en las soluciones de los problemas.

\section{Conceptos Clave de Cinemática}

\begin{description}

    \item[Posición:] Vector $\vec{r}(t)$ que especifica la ubicación de un punto.

    \item[Velocidad:] Primera derivada de la posición,
    $$\vec{v}(t) = \frac{d\vec{r}}{dt}$$
    Representa la tasa de cambio de posición.

    \item[Aceleración:] Segunda derivada de la posición (primera de la velocidad),
    $$\vec{a}(t) = \frac{d\vec{v}}{dt} = \frac{d^2\vec{r}}{dt^2}$$
    Representa la tasa de cambio de velocidad.

    \item[Velocidad Angular:] Para marcos móviles, $\vec{\omega}$ describe la rotación.

    \item[Aceleración Angular:] Primera derivada de la velocidad angular,
    $$\vec{\alpha} = \frac{d\vec{\omega}}{dt}$$

    \item[Velocidad Relativa:] En sistemas con dos marcos de referencia,
    $$\vec{v}_{P/B} = \vec{v}_P - \vec{v}_B$$

    \item[Composición de Velocidades:] Relaciona velocidades en marcos diferentes,
    $$\vec{v}_{P/A} = \vec{v}_{P/B} + \vec{v}_{B/A}$$

    \item[Restricciones:] Condiciones que relacionan las variables del sistema.
    Ejemplos: movimiento en una curva, rodadura sin deslizar, etc.

\end{description}

\section*{Objetivos del Capítulo}

Al completar este capítulo, deberás:
\begin{itemize}
    \item Usar competentemente los cuatro sistemas de coordenadas principales.
    \item Calcular velocidad y aceleración en cada sistema.
    \item Aplicar el Teorema del Transporte para marcos móviles.
    \item Analizar movimientos complejos usando cambios de marco.
    \item Resolver problemas con restricciones de movimiento.
\end{itemize}

\newpage
